\chapter{数据整理}\label{ch:g7:stats}

在进行任何计算之前，数据必须先经过收集、计数，
并加以展示。本章讨论如何把杂乱的观测结果整理为
清晰的表格或图形：频数、频率、分组数据、条形图与
扇形图。用一个数概括数据（如平均数、中位数）是
\cref{ch:g9:statproba}的内容。

\section{频数与频率}

\begin{definition}[频数与频率]\label{def:g7:stats:frequency}
对于调查中的每一个可能取值，其\emph{频数}\index{频数}是
该值出现的次数；其\emph{频率}\index{频率}
等于频数除以观测的总数：
\[
\text{频率} = \frac{\text{频数}}{\text{总数}} .
\]
频率可以写成分数、小数或百分数，
它们的和总是 $1$（即 $100\,\%$）。
\end{definition}

\begin{example}\label{ex:g7:stats:frequency}
某班 $25$ 名学生最喜欢的运动，从原始名单
整理为表格：
\begin{center}
\renewcommand{\arraystretch}{1.3}
\begin{tabular}{l|cccc|c}
运动 & 足球 & 游泳 & 网球 & 柔道 & 总计 \\
\hline
频数 & $10$ & $6$ & $5$ & $4$ & $25$ \\
频率 & $0.40$ & $0.24$ & $0.20$ & $0.16$ & $1$ \\
\end{tabular}
\end{center}
对于足球：$\frac{10}{25} = \frac{40}{100} = 0.40 = 40\,\%$。
频率一行的和必须为 $1$ --- 一项随时可用的免费检验。
\end{example}

\begin{remark}[为什么要用频率？]
频数不能直接用来比较规模不同的群体：$10$
名足球爱好者位于 $25$ 人的班级中（即 $40\,\%$），显得
比 $12$ 名身处 $120$ 人年级中的同学（即 $10\,\%$）更热情。
频率则能把所有数据放在同一尺度上。
\end{remark}

\section{分组整理数据}

\begin{example}[分组]\label{ex:g7:stats:classes}
$20$ 名学生的身高（单位：cm）在 $148$ 到 $172$ 之间。
若逐个列出每个身高，表格的列数将与学生人数几乎相同；
因此应分组，将数据分成若干\emph{组}：
\begin{center}
\renewcommand{\arraystretch}{1.3}
\begin{tabular}{l|ccc|c}
身高（cm） & $\intco{145}{155}$ & $\intco{155}{165}$ &
$\intco{165}{175}$ & 总计 \\
\hline
频数 & $5$ & $9$ & $6$ & $20$ \\
频率 & $0.25$ & $0.45$ & $0.30$ & $1$ \\
\end{tabular}
\end{center}
每个身高恰好计入一个组（$155$ 应归入
$\intco{155}{165}$ 中，而不是同时计入两组 --- 这正是区间取
半开形式的原因。分组会损失细节，但提高了可读性。
\end{example}

\section{图表}

\begin{method}[选择并绘制图表]\label{met:g7:stats:charts}
\begin{enumerate}
  \item \emph{条形图}：每个取值或每一组对应一个条形，条形的高度
        与频数（或频率）成正比例。适用于
        比较各个取值。
  \item \emph{扇形图}：每个取值对应一个扇形，其圆心角与
        频率成正比例 --- 整个圆（$360^\circ$）表示
        总数，因此
        \[
        \text{圆心角} = \text{频率} \times 360^\circ .
        \]
        适合表示各部分占整体的比例。
  \item 两种图表都应写明标题、标签，以及刻度或
        标注在图上的百分数。
\end{enumerate}
\end{method}

\begin{example}\label{ex:g7:stats:pie}
\cref{ex:g7:stats:frequency}中各项运动的圆心角：足球
$0.40 \times 360 = 144^\circ$；游泳 $0.24 \times 360 =
86.4^\circ$；网球 $72^\circ$；柔道 $57.6^\circ$。验算：
$144 + 86.4 + 72 + 57.6 = 360$。
\end{example}

\begin{omfigure}
\begin{tikzpicture}
\begin{axis}[
  width=6.6cm, height=5.2cm,
  ybar, bar width=14pt,
  symbolic x coords={football,swimming,tennis,judo},
  xtick=data,
  x tick label style={font=\footnotesize, rotate=25, anchor=east},
  ymin=0, ymax=12,
  ytick={2,4,6,8,10},
  ylabel={频数},
  label style={font=\small},
  tick label style={font=\small}]
  \addplot[fill=omDef!30, draw=omDef] coordinates
    {(football,10) (swimming,6) (tennis,5) (judo,4)};
\end{axis}
\end{tikzpicture}\hspace{0.8cm}%
\begin{tikzpicture}[scale=1.05]
  \fill[omDef!45] (0,0) -- (0:1.5) arc (0:144:1.5) -- cycle;
  \fill[omThm!35] (0,0) -- (144:1.5) arc (144:230.4:1.5) -- cycle;
  \fill[omMeth!35] (0,0) -- (230.4:1.5) arc (230.4:302.4:1.5) -- cycle;
  \fill[omProp!35] (0,0) -- (302.4:1.5) arc (302.4:360:1.5) -- cycle;
  \draw (0,0) circle (1.5);
  \foreach \a in {0,144,230.4,302.4} \draw (0,0) -- (\a:1.5);
  \node[font=\footnotesize] at (72:0.95) {足球};
  \node[font=\footnotesize] at (187:0.95) {游泳};
  \node[font=\footnotesize] at (266:0.95) {网球};
  \node[font=\footnotesize] at (331:0.95) {柔道};
\end{tikzpicture}

{\small 同一组数据既可用条形图表示（比较高度），
也可用扇形图表示（比较各部分在全班中所占的比例）。足球占
$144^\circ$，即占圆的 $40\,\%$。}
\end{omfigure}

\begin{example}[批判性地读图]\label{ex:g7:stats:critical}
纵轴从 $9$ 而非 $0$ 开始，会使一个
高度为 $10$ 的条形看上去比高度为 $9.1$ 的条形高十倍 --- 在视觉上
具有误导性，尽管数字本身是真实的。看任何图表时，首先应检查
纵轴从何处开始，以及每一格刻度代表多少
（见\cref{met:g6:propdata:read}）。
\end{example}

\section{练习}

\begin{exercise}[$\star$]\label{exo:g7:stats:1}
下面是 $20$ 名学生的成绩：$12$、$15$、$9$、$12$、$14$、
        $9$、$12$、$15$、$11$、$12$、$9$、$14$、$12$、$11$、$15$、$9$、$12$、
        $14$、$11$、$12$。编制频数表（每个成绩占一列）。
\end{exercise}

\begin{exercise}[$\star$]\label{exo:g7:stats:2}
将\cref{exo:g7:stats:1}中的频数表补充一行频率
（先以 $20$ 为分母写成分数，再写成百分数）。检验各频率之和。
\end{exercise}

\begin{exercise}[$\star$]\label{exo:g7:stats:3}
在一项对 $50$ 个家庭的调查中，$30$ 个家庭有一辆车。
“有一辆车”的频率是多少？分别用小数和百分数表示。
在另一项调查中，$45$ 个家庭（共 $90$ 个家庭）有一辆车，
哪一项调查显示的比例更高？
\end{exercise}

\begin{exercise}[$\star$]\label{exo:g7:stats:4}
下面是 $15$ 只狗的质量（千克）：$8$、$23$、$31$、$12$、$9$、$27$、$35$、
$14$、$22$、$18$、$29$、$11$、$25$、$33$、$16$。将它们分入
下列组段 $\intco{5}{15}$、$\intco{15}{25}$、$\intco{25}{35}$、
$\intco{35}{45}$ 中，并给出各组的频数。
\end{exercise}

\begin{exercise}[$\star$]\label{exo:g7:stats:5}
根据下表画出条形图：
\begin{center}
\renewcommand{\arraystretch}{1.2}
\begin{tabular}{l|cccc}
宠物数 & $0$ & $1$ & $2$ & $3$ \\
\hline
家庭数 & $8$ & $12$ & $6$ & $4$ \\
\end{tabular}
\end{center}
\end{exercise}

\begin{exercise}[$\star$]\label{exo:g7:stats:6}
对\cref{exo:g7:stats:5}中的表格，求各数值的频率及
对应扇形图的圆心角（检验：总和应为 $360^\circ$）。
\end{exercise}

\begin{exercise}[$\star$]\label{exo:g7:stats:7}
一幅关于最喜爱季节的扇形图显示：夏季 $162^\circ$、春季
$90^\circ$、秋季 $54^\circ$、冬季 $54^\circ$。各季节分别占百分之几？
若有 $60$ 人作答，其中选择夏季的有多少人？
\end{exercise}

\begin{exercise}[$\star\star$]\label{exo:g7:stats:8}
在 A 班，$12$ 名学生（共 $30$ 名）步行上学；在 B 班，
$14$ 名学生（共 $40$ 名）。哪个班步行上学的比例更高？
请用频率而非频数来说明理由。
\end{exercise}

\begin{exercise}[$\star\star$]\label{exo:g7:stats:9}
某杂志刊登了一幅月销售额条形图：$9\,800$，接着是
$10\,000$，再是 $10\,100$，纵轴从 $9\,700$ 开始。
请描述读者所看到的情况，以及若纵轴从 $0$ 开始，
则会呈现怎样的不同结果？
\end{exercise}

\begin{exercise}[$\star\star$]\label{exo:g7:stats:10}
某频数表包含三个数值，其频率分别为 $0.35$、$0.4$ 和
$f$。求 $f$。若总频数为 $80$，求出这三个频数。
\end{exercise}

\begin{exercise}[$\star\star\star$]\label{exo:g7:stats:11}
某校学生中，$55\,\%$ 为女生。在女生中，$40\,\%$ 在食堂就餐；
在男生中，$60\,\%$ 在食堂就餐。若全校共
$400$ 名学生，在食堂就餐的有多少人？（请逐步算出四个
群体的人数。）这相当于总体的百分之几？
\end{exercise}

\section{问题：数据如何误导人 --- 以及如何识破}

\begin{problem}[{周末问题 --- 频数与频率之争，
让一个国家上当的民意调查，以及
两家医院的悖论}]
\label{pb:g7:stats:1}
数字本身不会说谎，但它们可能被用来误导：巨大的频数可能掩盖
很小的频率，规模庞大的调查也可能毫无价值，
而且 --- 最奇怪的是 --- 一家医院可能在
\emph{每一类}患者上都胜过对手，却仍在
总体数字上落败。本题将逐一揭示这三个陷阱，
仅借助本章的频数与频率，
（见\cref{def:g7:stats:frequency}）。

\medskip
\noindent\textbf{第一部分 --- 频数并非频率。}
\begin{enumerate}
  \item A 校回收了 $120$ 个果汁盒，共 $400$ 个；
        B 校回收了 $90$ 个果汁盒，共 $250$ 个。哪所学校回收的
        \emph{频数}更多？
        哪所学校的\emph{频率}更高？哪所学校应获得
        回收奖？
  \item 对一份新课表进行调查：$18$ 名（共 $30$ 名）教师喜欢它，
        另有 $210$ 名（共 $600$ 名）学生也喜欢。计算
        两个频率，再计算 $630$ 人合在一起的
        频率。为什么合并后的结果会
        与学生的频率如此接近？
  \item 一幅已发表的扇形图中，各扇形标注的数值分别为 $45\,\%$、
        $30\,\%$、$20\,\%$ 和 $10\,\%$。在没有其他信息的情况下，
        你如何断定其中出了错？
  \item 一张基于 $40$ 次观测的频数表残缺不全：
        频数分别为 $14$、$10$、$?$、$?$，频率分别为 $0.35$、
        $0.25$、$0.15$、$?$。请补全缺失的数据。
  \item 一项调查的频率最终为 $0.40$、$0.35$ 和
        $0.25$。计算三个扇形图的圆心角
        （见\cref{met:g7:stats:charts}），并检验它们是否恰好组成
        一个完整的圆。
\end{enumerate}

\medskip
\noindent\textbf{第二部分 --- 骗过一个国家的民调。}
1936 年，一家美国杂志向
从电话簿和汽车登记名单中取得的地址寄出一千万张选票，
收到 $2.4$ 百万份回复，并预测兰登将获得
压倒性胜利。一位年轻的统计学家乔治·盖洛普
只调查了约五万人，这些人经过挑选，
以与全体人口相仿，并预测了相反的结果。
罗斯福大获全胜。
\begin{enumerate}[resume]
  \item 寄出的选票中，收回的频率是多少？
  \item 1936 年，电话和汽车都还是奢侈品。请用一两句话解释，
        这家杂志的样本虽然规模庞大，却为何注定失败——
        并说明第 1 题中的哪一条教训
        在全国范围内再次出现。
  \item 一个简化模型。某镇有 $1\,000$ 名富裕选民，其中
        $70\,\%$ 支持 L，另有 $9\,000$ 名普通选民，其中
        $30\,\%$ 支持 L。计算 L 在该镇的
        真实支持率。“电话簿民调”覆盖了 $500$ 名富裕选民和
        $500$ 名普通选民：它预测的支持率是多少？
  \item 无声的陷阱：某公司有 $200$ 名不满意的客户，
        以及 $800$ 名满意的客户。公司对全体客户进行调查，
        不满意者中有 $40\,\%$ 作出答复，满意者中只有 $10\,\%$ 作出答复。
        计算两类答复各有多少份，以及
        不满意者在\emph{全部答复之中}的频率，
        并与真实频率作比较。
  \item 修正第 8 题中的民调：保持 $1\,000$
        次访问不变，但按该镇的真实比例进行分配。
        修正后的民调预测了什么结果？
\end{enumerate}

\medskip
\noindent\textbf{第三部分 --- 两家医院的
悖论。}
两家医院公布了各自按病情轻重分类的治愈人数：
\begin{center}
\renewcommand{\arraystretch}{1.3}
\begin{tabular}{l|cc}
 & 轻症 & 重症 \\
\hline
医院 A & $90$ 人治愈（共 $100$ 人）& $280$ 人治愈（共 $400$ 人）\\
医院 B & $340$ 人治愈（共 $400$ 人）& $65$ 人治愈（共 $100$ 人）\\
\end{tabular}
\end{center}
\begin{enumerate}[resume]
  \item 分别计算医院 A 轻症患者的治愈频率、
        重症患者的治愈频率，以及总体治愈频率（全部 $500$ 名患者）。
  \item 对医院 B 作同样的三项计算。哪家医院
        在轻症中胜出？在重症中胜出？
  \item 比较两家医院的总体治愈频率，直截了当地指出
        所出现的奇怪现象。
  \item 解释其中的诀窍：比较两家医院患者的
        \emph{构成}，并说明为什么总体频率
        可能与两类频率的结论都相反。若你
        是重症患者，应选择哪家医院？
  \item 一则报纸标题写道：“医院 B 的治愈率更高，
        即 $81\,\%$ 对 $74\,\%$。”
        请写出本题让你学会撰写的两句话致编辑信——
        第一句说明标题忽略了什么，
        第二句说明一般的道理
        （比较同类对象，并注意权重）。
\end{enumerate}
\end{problem}
